% ЗАКЛЮЧЕНИЕ, без номер.
\section*{ЗАКЛЮЧЕНИЕ}
\label{sec:conclusion}
\addcontentsline{toc}{section}{ЗАКЛЮЧЕНИЕ}

В рамките на проекта е проектирана и реализирана паралелна рамка за метаевристична
оптимизация на комбинаторни задачи, в която задачата, метаевристиката и схемата за
паралелно изпълнение са три независими точки на разширение. Реализирани са четири
метаевристики и три схеми за паралелно изпълнение, а към тях е изградена методика за
измерване, която фиксира входните инстанции, критерия за спиране, броя повторения и
начина на обобщаване. Рамката се изгражда без нито една външна зависимост и се покрива
от 28 автоматични теста.

Измерванията дават три резултата, които си заслужава да бъдат отделени. Първо,
времето до фиксиран праг качество намалява до 7,20 пъти при дванадесет нишки за
независимия многократен старт, като при две до четири нишки ускорението е свръхлинейно,
до 5,43 при четири. Обяснението не е в машината, а в разпределението: междуквартилният
размах при една нишка е 84,22 ms при медиана 87,45 ms, тоест опашката е тежка и
минимумът от няколко независими тегления пада по-бързо от $1/p$. Второ, кооперативната
схема подобрява медианното качество с 0,93 на сто спрямо независимата при задачата на
търговския пътник и я влошава при раницата, тоест обменът не е универсално подобрение;
цената му е около 1 на сто от процесорното време. Трето, подредбата на четирите
метаевристики зависи от задачата и не се пренася между задачите.

\subsection*{Предимства и недостатъци на решението}

Предимствата се формулират спрямо измереното. Разделянето на договорите работи:
добавянето на трета задача не наложи промяна в нито един от четирите алгоритъма, а
смяната на схемата за изпълнение не наложи промяна в нито една задача. Единната мярка
за бюджет в оценявания на целевата функция позволява алгоритми, чиито итерации се
различават с четири порядъка по цена, да бъдат сравнявани и да обменят решения
смислено. Инкременталното оценяване се изплаща и е измерено колко: при 1000 града
пълното преоценяване на всяка итерация би струвало 13 пъти повече от цялата итерация.

Недостатъците са три и всеки от тях е показан от измерване, а не от преглед на кода.
Първо, качеството зависи от съседството поне толкова, колкото от избраната стратегия:
едно и също табу търсене стига до 27579 или до 24965 в зависимост от това дали
обхватът на хода е ограничен, тоест сравнението между метаевристики е валидно само при
еднакво добре подбрано съседство и за конкретните настройки. Второ, брояча за време в
синхронизация мери различни неща при кооперативната и при островната схема, тоест двете
колони не са пряко сравними; това е дефект на инструментирането, признат в
§\ref{subsec:res-cooperation}. Трето, единадесет повторения не разграничават съседни
конфигурации, когато междуквартилният размах е съизмерим с медианата, което личи по
отделните точки в Табл.~\ref{tab:speedup}, отклоняващи се от тенденцията.

\subsection*{Насоки за по-нататъшна работа}

Работата може да бъде продължена в няколко посоки. Най-близката следва от отворения
въпрос в §\ref{subsec:res-summary}: разграничаването между споделяне на изпълнителни
единици и ограничение по памет като причина за сплескването на ускорението изисква
хардуерни броячи и следователно среда, в която такива могат да бъдат снети. Втората е
разпределеното изпълнение върху няколко машини, което поставя въпроса за обмена при
съществено по-висока цена на комуникацията, тоест променя баланса, измерен тук. Третата
е автоматичната настройка на параметрите, която би премахнала едно от ограниченията,
признати в §\ref{subsec:method-threats}, а именно че сравнението е между конкретни
настройки, а не между методите изобщо. Задачно ориентираната среда за изпълнение остава
разгледана, но неизмерена алтернатива по §2.2.
